%% Copyright 2026 Petr Socha
%%
%% This work may be distributed and/or modified under the
%% conditions of the LaTeX Project Public License, either version 1.3c
%% of this license or (at your option) any later version.
%% The latest version of this license is in
%%   https://www.latex-project.org/lppl.txt
%% and version 1.3c is part of all distributions of LaTeX
%% version 2008 or later.
%%
%% This work has the LPPL maintenance status `maintained'.
%% The Current Maintainer of this work is Petr Socha.
%% Maintenance and bug reports: https://github.com/petrsocha/pslib-latex/issues
%%
%% The components of this work are listed in MANIFEST.txt.
%% The school logo is a separately licensed third-party asset.

% ============================================================================
% REALIZACE -- ZDROJOVÝ KÓD A SPECIÁLNÍ ZNAKY
% ============================================================================
\chapter{Realizace}
\label{chap:realizace}

\section{Výpis zdrojového kódu}

% Krátký zdrojový kód vložte do prostředí lstlisting.
% Výpis dostane automatické číslo, popisek a volitelně čísla řádků.
% Diakritika v komentářích a řetězcích Pythonu je podporována.
\begin{lstlisting}[language=Python,caption={Výpočet průměru naměřených hodnot},label={lst:prumer}]
def prumer(hodnoty):
    if not hodnoty:
        return None
    return sum(hodnoty) / len(hodnoty)

# Český komentář i řetězec ve zdrojovém kódu
print("Průměrná hodnota:", prumer([20, 22, 24]))
\end{lstlisting}

Jednoduchou funkci v~jazyce Python ukazuje \autoref{lst:prumer}.
U~rozsáhlejších programů je vhodné uvést jen důležité úseky a~úplné
zdrojové soubory přiložit samostatně.

\section{Speciální znaky}

% Znaky %, _, &, # a $ maji v LaTeXu speciální význam.
% V běžném textu se zapisují jako \%, \_, \&, \# a \$.
Cesta k~souboru může vypadat jako \texttt{data\_2026.csv}.
Hodnota 50\,\% znamená polovinu celku; text \texttt{A\&B} obsahuje znak
ampersandu. Příkaz \verb|\texttt{...}| sází neproporcionálním písmem.
